% Appendix Q: Quarantine
% External data fitting procedures
% Status: FILLED with calibration details

\chapter{Quarantine: Fitting Procedures}
\label{app:quarantine}

\begin{tcolorbox}[colback=red!5,colframe=red!70!black,title=Warning]
This appendix contains fitting procedures and external data comparisons.
Content here is Layer B (calibration) and does not affect Layer A (derivations).
All fitted parameters are marked with [Cal] when used in main text.
\end{tcolorbox}

% ============================================================================
\section{Baseline Lane Coefficients}
\label{sec:q:gn}
% ============================================================================

The baseline lane coefficients are determined by regression on empirical
closed-4 release data:
\begin{equation}
    \log_{10}(t_{1/2}) = a \times X + b
\end{equation}
where $X = Z_{\text{eff}}/\sqrt{E_{\text{rel}}}$ is the barrier-to-release ratio.

\subsection{Fitted Values}

\begin{center}
\begin{tabular}{lcc}
    \toprule
    Coefficient & Value & Uncertainty \\
    \midrule
    $a$ & $1.632$ & $\pm 0.028$ \\
    $b$ & $-50.75$ & $\pm 0.91$ \\
    \bottomrule
\end{tabular}
\end{center}

\subsection{Fitting Procedure}

\begin{enumerate}
    \item Dataset: 106 closed-4 emitters with $A > 150$
    \item Regression: Ordinary least squares on $\log_{10}(t_{1/2})$ vs.\ $X$
    \item Outlier handling: None (all data points included)
    \item Goodness of fit: $R^2 = 0.98$
\end{enumerate}

\subsection{Data Sources}

Empirical release times from published measurements (Layer B only).
Specific dataset identifiers suppressed from Layer A.

% ============================================================================
\section{Prefactor $p$ in Coordination Function}
\label{sec:q:prefactor}
% ============================================================================

The prefactor $p$ in $n(A) = p \times A^{1/3}$ is calibrated by anchoring
$n(208) \approx 36$ for the heaviest doubly-stable cluster.

\subsection{Calibration}

\begin{align}
    n(208) &= p \times 208^{1/3} = p \times 5.925 \\
    36 &= p \times 5.925 \\
    p &= 6.08 \approx 6.1
\end{align}

\subsection{Sensitivity Analysis}

\begin{center}
\begin{tabular}{cccc}
    \toprule
    Prefactor $p$ & Mean $|\Delta|$ (dex) & Pass rate & Status \\
    \midrule
    6.00 & 0.93 & 5/6 & Acceptable \\
    6.05 & 0.89 & 5/6 & Acceptable \\
    \textbf{6.10} & \textbf{0.88} & \textbf{5/6} & \textbf{Selected} \\
    6.15 & 0.91 & 5/6 & Acceptable \\
    6.20 & 0.95 & 5/6 & Acceptable \\
    \bottomrule
\end{tabular}
\end{center}

The model is robust to prefactor variations within $\pm 0.1$.

% ============================================================================
\section{Frustration Coefficient $g$}
\label{sec:q:g}
% ============================================================================

The coefficient $g$ in the frustration correction $\Delta \approx g \cdot d(n)$
is determined by regression on baseline residuals.

\subsection{Fitted Value}

\begin{center}
\begin{tabular}{lcc}
    \toprule
    Coefficient & Value & Uncertainty \\
    \midrule
    $g$ & $-1.763$ & $\pm 0.142$ \\
    \bottomrule
\end{tabular}
\end{center}

\subsection{Sign Interpretation}

The negative sign $g < 0$ indicates that:
\begin{itemize}
    \item Higher coordination frustration $\to$ faster release
    \item Frustrated configurations have reduced effective barriers
    \item Consistent with enhanced preformation or reduced binding
\end{itemize}

\subsection{Robustness Tests}

\begin{center}
\begin{tabular}{lcc}
    \toprule
    Control Variable & $g$ Value & Change \\
    \midrule
    None (baseline) & $-1.64$ & --- \\
    + Boundary proxy & $-1.46$ & $-11\%$ \\
    + Preformation proxy & $-1.53$ & $-7\%$ \\
    + Both proxies & $-1.71$ & $+4\%$ \\
    \bottomrule
\end{tabular}
\end{center}

The coefficient $g$ is stable across different model specifications.

% ============================================================================
\section{High-Coordination Benchmark Dataset}
\label{sec:q:benchmark}
% ============================================================================

This section documents the benchmark dataset used in Chapter~\ref{ch:high_coord}.

\subsection{Case Definitions}

\begin{center}
\begin{tabular}{lccccc}
    \toprule
    Case ID & $A$ & $X$ & $E_{\text{rel}}$ (MeV) & Status \\
    \midrule
    C-1 & 289 & 35.79 & 9.82 & Observed \\
    C-2 & 290 & 37.00 & 9.19 & Observed (marginal) \\
    C-3 & 290 & 35.08 & 10.41 & Observed \\
    C-4 & 293 & 34.45 & 10.67 & Observed \\
    C-5 & 294 & 34.03 & 10.81 & Observed \\
    C-6 & 294 & 32.24 & 11.65 & Observed \\
    C-7 & 298 & 34.50 & 10.50 & Predicted \\
    C-8 & 302 & 35.50 & 10.00 & Predicted \\
    C-9 & 304 & 36.38 & 9.50 & Predicted \\
    \bottomrule
\end{tabular}
\end{center}

\subsection{Selection Criteria}

\begin{enumerate}
    \item Mass index $A \geq 289$ (high-coordination regime)
    \item Coordination function $n(A) > 40$ (forbidden zone)
    \item Primary closed-4 emission channel confirmed
    \item Release time measurement available (for observed cases)
\end{enumerate}

\subsection{Anti-Tuning Statement}

The benchmark dataset was selected \textbf{before} parameter optimization.
No cases were excluded based on model performance. The coefficient $g$ was
determined on a separate training set ($A < 252$) and applied without
adjustment to the high-coordination benchmark.

% ============================================================================
\section{Performance Metrics Computation}
\label{sec:q:metrics}
% ============================================================================

\subsection{Definitions}

\begin{itemize}
    \item \textbf{Baseline residual:}
          $\Delta := \log_{10}(t_{\text{obs}}) - \log_{10}(t_{\text{BL}})$
    \item \textbf{Corrected residual:}
          $\Delta_{\text{corr}} := \log_{10}(t_{\text{obs}}) - \log_{10}(t_{\text{pred}})$
    \item \textbf{Mean absolute error:}
          $\text{MAE} = \frac{1}{N}\sum_i |\Delta_i|$
    \item \textbf{Pass rate:}
          Fraction of cases with $|\Delta| < 1.5$ dex
\end{itemize}

\subsection{Results Summary}

On the high-coordination benchmark (Cases C-1 through C-6):

\begin{center}
\begin{tabular}{lcc}
    \toprule
    Metric & Baseline & Corrected \\
    \midrule
    Mean $|\Delta|$ & 6.16 dex & 0.88 dex \\
    Median $|\Delta|$ & 6.02 dex & 0.86 dex \\
    Max $|\Delta|$ & 8.15 dex & 1.40 dex \\
    Pass rate & 0/6 & 5/6 \\
    \bottomrule
\end{tabular}
\end{center}

% ============================================================================
\section{Cross-Validation Results}
\label{sec:q:cv}
% ============================================================================

Five-fold cross-validation on the full dataset ($N = 106$):

\begin{center}
\begin{tabular}{ccc}
    \toprule
    Fold & Train $R^2$ & Test MAE (dex) \\
    \midrule
    1 & 0.979 & 0.42 \\
    2 & 0.981 & 0.38 \\
    3 & 0.978 & 0.51 \\
    4 & 0.980 & 0.44 \\
    5 & 0.977 & 0.47 \\
    \midrule
    Mean & 0.979 & 0.44 \\
    \bottomrule
\end{tabular}
\end{center}

Cross-validation confirms model stability across data splits.

% ============================================================================
\section{Declaration}
\label{sec:q:declaration}
% ============================================================================

All fitting procedures are documented in this appendix. The main text
(Layer A) contains no fitted parameters except where explicitly marked
with [Cal]. Baseline empirical values are marked [BL].

\textbf{Guardrails:}
\begin{itemize}
    \item No post-hoc case exclusion based on model performance
    \item Prefactor $p$ anchored to physical constraint ($n(208) = 36$)
    \item Coefficient $g$ determined on training set, applied without
          adjustment to test set
    \item All sensitivity analyses documented
\end{itemize}

